\section{Introduction}
\label{sec:intro}

3D semantic understanding is a fundamental capability in computer vision that enables machines to interpret and reason about the 3D world, with practical use in robotics~\cite{rosinol20203ddynamicscenegraphs}, autonomous navigation~\cite{zhou2020cylinder3d, tian2024drivevlmconvergenceautonomousdriving}, augmented reality, and interactive systems. While traditional 3D scene understanding~\cite{gupta2015indoor} focuses on extracting geometric structure and predefined information from visual data, newer approaches integrate natural language. This emerging field, known as 3D language understanding~\cite{peng2023openscene3dsceneunderstanding, kerr2023lerf, gu2023conceptgraphsopenvocabulary3dscene, liu2023weakly, ha2022semantic, zhang2023clip, takmaz2023openmask3d, ding2023pla}, aims to establish rich and meaningful connections between linguistic expressions and 3D visual representations such as point clouds, meshes, and RGB-D images.

Recent breakthroughs in 3D scene representation~\cite{nerf, kerbl3Dgaussians, sun2022direct, yu2021plenoxels}, particularly 3D Gaussian Splatting~\cite{kerbl3Dgaussians, Yu2024MipSplatting}, have revolutionized scene reconstruction and rendering by offering an efficient and high-quality approach to modeling 3D environments. While 3D Gaussian Splatting excels at geometric reconstruction and novel view synthesis, its potential for semantic understanding remains largely unexplored. Recent research has begun addressing this limitation by incorporating language features into the Gaussian representations, enabling open-vocabulary scene understanding through natural language queries. 

Vision-language models like CLIP~\cite{caron2021emerging} and LSeg~\cite{ghiasi2022scaling} have demonstrated the power of joint visual-linguistic understanding by training on vast image-text pairs, while foundation models like Segment Anything (SAM)~\cite{kirillov2023segment} and DINO-V2~\cite{oquab2024dinov2} have shown the potential of large-scale visual learning.
Although these models have changed 2D image understanding completely, extending their capabilities to 3D domains remains challenging~\cite{qin2024langsplat, peng20243d, wu2024opengaussian, shi2024language, ye2025gaussian, bhalgat2024n2f2}: directly learning 3D-language representations lacks paired data and 2D vision-language models typically produce expressive high-dimensional features, which poses an often infeasible computational cost when integrated in complex 3D scenes.

Current approaches address this challenge through scene-specific feature autoencoders~\cite{qin2024langsplat} or quantization to create a compact feature space~\cite{shi2024language}. However, these scene-specific compression approaches fundamentally limit scalability and generalization. While effective for individual environments, they create barriers for multi-scene applications by requiring separate training for each scene, increasing both computational and storage overhead. Moreover, these scene-specific solutions make it difficult to transfer knowledge between environments and limit the potential for open-world understanding.

We take a fundamentally different approach by proposing a training-free method, that avoids repeated iteration over views, as well as loss computation in the high-dimensional feature space. Overall, our contribution can be summarized as follows:
\begin{itemize}
    \item A training-free global optimization approach for Language 3D Gaussian Spatting, improving speed by two orders of magnitude.
    \item A reasoning by rendering approach to avoiding expensive feature rendering by following a probabilistic formulation of the Gaussian Splatting forward process.
    \item Compatibility with uncompressed and arbitrary dimensional language features.
\end{itemize}

\begin{table}[h]
    \centering
    \footnotesize
    \begin{tabular}{@{} lcccc|cc}
        \toprule
        & \multicolumn{4}{c}{Requires} & & \\ % Common heading spanning columns 2 to 5
        %\cmidrule(lr){2-5} % Underline the 'Requires' heading from column 2 to 5
        Method
        & \rotatebox{90}{\parbox{1.8cm}{\raggedright Feature Compression}}
        & \rotatebox{90}{\parbox{1.8cm}{\raggedright Spatial Compression}}
        & \rotatebox{90}{\parbox{1.8cm}{\raggedright Back Propagation}}
        & \rotatebox{90}{\parbox{1.8cm}{\raggedright Scene Specific Model}}
        & \rotatebox{90}{\parbox{1.8cm}{\raggedright Probabilistically Grounded}}        
        & \rotatebox{90}{\parbox{1.8cm}{\raggedright Fast}} \\
        \midrule
        LERF~\cite{kerr2023lerf} & \xmark & \xmark & \cmark & \cmark & \xmark & \xmark \\
        GS-Grouping~\cite{ye2025gaussian} & \cmark & \xmark & \cmark & \cmark & \xmark & \xmark \\
        Feature-3DGS~\cite{zhou2024feature} & \cmark & \xmark & \cmark & \cmark & \xmark & \xmark \\
        LEGaussians~\cite{shi2024language} & \cmark & \xmark & \cmark & \cmark & \xmark & \xmark \\
        GOI~\cite{qu2024goi3dgaussiansoptimizable} & \cmark & \xmark & \cmark & \cmark & \xmark & \xmark \\
        LangSplat~\cite{qin2024langsplat} & \cmark & \xmark & \cmark & \cmark & \xmark & \xmark \\
        FMGS~\cite{zuo2024fmgs} & \xmark & \cmark & \cmark & \cmark & \xmark & \xmark \\
        Occam's LGS & \xmark & \xmark & \xmark & \xmark & \cmark & \cmark \\
        \bottomrule
    \end{tabular}
    \caption{Our approach challenges the necessity of complex modules to lift 2D language features to Gaussian Splats, resulting in a simple and fast, yet probabilistically grounded method.}
    \vspace{-5mm}
    \label{tab:method-comparison}
\end{table}